\label{chap.2}
\label{chap:literature}

\section{Introduction}

Lithium-ion batteries are widely used in electric vehicles because of
their high energy density, high efficiency, and long cycle life.
However, safe and efficient battery operation requires effective
battery management systems (BMS). A BMS is responsible for monitoring
battery states, controlling charging and discharging, estimating
battery health, and preventing unsafe operating conditions.

Conventional battery charging methods are generally based on predefined
rules, while model-based methods use mathematical representations of
battery behaviour. More recently, reinforcement learning (RL) and
multi-agent reinforcement learning (MARL) have been investigated for
adaptive battery control. These methods provide opportunities to
optimize charging performance while considering battery health and
operating constraints.

This chapter reviews the major approaches related to battery charging,
battery state estimation, reinforcement learning, multi-agent battery
management, and safe reinforcement learning. The limitations of these
approaches are then analysed to identify the research gaps addressed by
the proposed multi-agent battery management framework.

\section{Conventional Battery Charging Methods}

The constant-current constant-voltage (CC-CV) method is one of the
most widely used charging strategies for lithium-ion batteries. During
the constant-current stage, the battery is charged using a predefined
current until the terminal voltage reaches a specified limit. The
charger then enters the constant-voltage stage, where the voltage is
maintained while the charging current gradually decreases.

Other charging strategies, including multi-stage constant-current and
pulse charging, have also been investigated to improve charging
performance. These approaches can modify the charging current according
to predefined operating conditions.

However, conventional charging strategies generally use fixed rules.
They do not dynamically adapt the charging policy according to changes
in battery ageing, internal resistance, temperature, or operating
conditions. Consequently, a charging strategy that performs well for a
new battery may not provide the same performance for an aged battery.

Therefore, conventional charging methods provide a reliable baseline,
but their limited adaptability creates a need for more intelligent
charging strategies.

\section{Battery Management Systems and State Estimation}

A battery management system continuously monitors important battery
variables such as voltage, current, temperature, state of charge (SoC),
and state of health (SoH). Accurate estimation of these states is
essential because they directly influence charging decisions and safety.

State estimation methods can generally be divided into model-based and
data-driven approaches. Model-based methods use battery models together
with measurements to estimate internal battery states. The Extended
Kalman Filter (EKF) is a commonly used approach for estimating SoC and
other battery parameters.

Battery ageing introduces additional difficulties because parameters
such as internal resistance and available capacity change over the
battery lifetime. An estimator designed using nominal battery
parameters may therefore become less accurate when the battery
degrades.

This indicates the importance of considering battery ageing when
designing both state estimation and charging-control strategies.

\section{Battery Ageing and State-of-Health Estimation}

Battery ageing is associated with capacity loss, increasing internal
resistance, and changes in the electrochemical behaviour of the cell.
These changes affect the voltage, temperature, charging time, and
available energy of the battery.

State-of-health estimation methods attempt to quantify the degradation
of a battery using measurements such as voltage, current, temperature,
capacity, and internal resistance. Both model-based and
machine-learning-based approaches have been investigated for this
purpose.

Accurate SoH information is particularly important for charging
control. A charging current that is acceptable for a healthy battery
may become unsuitable as internal resistance increases with ageing.

However, many battery-control approaches are evaluated using nominal
or relatively fixed battery parameters. This creates a gap between
battery-health estimation and adaptive charging control.

\section{Model-Based Battery Charging Control}

Model Predictive Control (MPC) is an important model-based approach for
battery charging. MPC predicts the future behaviour of the battery over
a finite time horizon and determines a control action while considering
system constraints.

Electrochemical models such as the Single Particle Model (SPM) and
Doyle--Fuller--Newman (DFN) models can be used to represent battery
dynamics. These models can provide information about internal battery
states and can be incorporated into constrained charging control.

The main advantage of MPC is its ability to explicitly consider
constraints such as voltage, temperature, and charging-current limits.
However, its performance depends strongly on the accuracy of the
battery model and its parameters. Detailed electrochemical models can
also introduce considerable computational requirements.

Furthermore, battery parameters change with ageing. Therefore, a
controller based on fixed model parameters may not accurately represent
an aged battery.

Thus, model-based control provides strong constraint-handling
capabilities but remains sensitive to model accuracy and parameter
variation.

\section{Reinforcement Learning for Battery Charging}

Reinforcement learning provides an alternative approach in which an
agent learns a charging policy through interactions with a battery
environment. Deep reinforcement learning methods such as Deep
Deterministic Policy Gradient (DDPG), Deep Q-Networks (DQN), and
Proximal Policy Optimization (PPO) can be used for continuous or
discrete charging-control problems.

An RL charging agent typically receives battery states such as SoC,
voltage, and temperature and produces a charging action. The agent is
then evaluated using a reward function.

RL provides an important advantage because the controller can learn
nonlinear relationships between battery states and charging actions.
However, several objectives may need to be considered simultaneously.
For example, a charging controller may need to minimize charging time
while also reducing temperature rise and battery degradation.

These objectives are commonly combined into a single scalar reward.
The resulting behaviour depends on the selected reward weights. A
change in these weights can therefore change the balance between
charging speed, battery health, and safety.

Another limitation is that a learned policy does not inherently
guarantee that hard physical constraints will always be satisfied.

\section{Multi-Agent Reinforcement Learning for Battery Management}

Multi-Agent Reinforcement Learning (MARL) extends reinforcement
learning by allowing multiple agents to interact with the same
environment. MARL has been investigated in areas such as energy
management, energy-storage systems, and battery systems.

In many existing battery-related MARL approaches, agents are assigned
according to physical components. For example, different cells or
energy-storage components may be represented by individual agents.
These agents then cooperate to achieve a common objective.

An alternative approach is to divide the control problem according to
functional objectives rather than physical components. For battery
management, different agents can therefore be responsible for
different aspects of the decision-making process.

Such functional decomposition can separate charging objectives from
battery-health considerations and provide a more modular control
architecture.

However, the literature provides limited investigation of functional
multi-agent decomposition specifically for EV battery charging while
simultaneously considering battery health and safety constraints.

\section{Safe and Constraint-Aware Reinforcement Learning}

Safety is a critical requirement in battery charging because excessive
voltage, temperature, or current can damage the battery.

Safe reinforcement learning attempts to incorporate safety constraints
into the learning process. Approaches include constrained policy
optimization, safety penalties, safety filters, and control-barrier
function-based methods.

A major limitation of penalty-based approaches is that safety is
represented indirectly through the reward. A sufficiently large
penalty may discourage unsafe behaviour, but it does not necessarily
provide a deterministic guarantee that a physical constraint will
never be violated.

Safety-filter approaches can provide an additional layer between the
learned policy and the physical system. However, their effectiveness
depends on the accuracy of the observed battery state and the model
used to determine safe actions.

For battery systems, this motivates an architecture in which the
learning agents propose actions while an independent safety mechanism
enforces hard operating limits.

\section{Empirical Battery Data and Evaluation}

Battery-control research commonly uses battery simulators and
experimental datasets to evaluate algorithms. Battery ageing datasets
provide measurements of voltage, current, temperature, capacity, and
other variables across repeated charge and discharge cycles.

However, differences in data collection procedures, operating
conditions, battery ageing behaviour, and failure definitions can
affect the interpretation of experimental results.

Similarly, remaining useful life (RUL) prediction requires careful
consideration of battery degradation and available failure data.
Evaluation becomes difficult when only a limited number of batteries
reach the defined end-of-life condition.

Therefore, battery-control research requires both appropriate physical
modelling and careful evaluation of the underlying experimental data.

\section{Summary of the Reviewed Works}

The reviewed approaches can be summarized according to their primary
strengths and limitations.

\begin{table}[htbp]
\centering
\caption{Summary of reviewed battery management approaches.}
\label{tab:literature_summary}
\begin{tabular}{p{3.0cm}p{4.0cm}p{4.5cm}}
\toprule
\textbf{Approach} & \textbf{Main Strength} & \textbf{Main Limitation} \\
\midrule
CC-CV and rule-based control &
Simple and reliable implementation &
Limited adaptation to ageing and changing conditions \\

MPC &
Explicit handling of system constraints &
Dependent on model accuracy and computational resources \\

State estimation &
Provides estimates of SoC, SoH, and battery parameters &
Accuracy can degrade when battery parameters change with ageing \\

Single-Agent RL &
Adaptive nonlinear control and optimization &
Multiple objectives are often combined into a single reward \\

Multi-Agent RL &
Allows decomposition of complex control problems &
Limited work on functional decomposition for battery charging \\

Safe RL &
Introduces safety constraints into learned control &
Safety may depend on reward design, observations, or model accuracy \\
\bottomrule
\end{tabular}
\end{table}

The reviewed literature shows that no single approach simultaneously
provides adaptive charging control, explicit consideration of battery
health, functional decomposition of control objectives, and an
independent mechanism for enforcing safety constraints.

\section{Research Gap}

Based on the reviewed literature, the following research gaps are
identified.

\subsection{Gap 1: Limited Ageing-Aware Charging Control}

Many charging-control approaches use fixed or nominal battery
parameters. Since battery resistance and capacity change with ageing,
a controller designed for a healthy battery may not remain optimal
throughout the battery lifetime.

\subsection{Gap 2: Objective Entanglement in Single-Agent RL}

Single-agent reinforcement learning commonly combines charging time,
battery degradation, temperature, and safety into one reward function.
The resulting behaviour can become sensitive to reward-weight
selection.

A more modular approach is required to separate different functional
objectives.

\subsection{Gap 3: Lack of Independent Hard Safety Enforcement}

Reward penalties alone cannot guarantee that a learned policy will
always satisfy physical voltage, temperature, and current constraints.

An independent safety mechanism is therefore required to verify and,
when necessary, modify the action proposed by the learning agent.

\subsection{Gap 4: Limited Integration of Battery Health with Charging Control}

Battery-health estimation and charging control are often treated as
separate problems. There is a need to incorporate health information
directly into the charging decision process so that charging
performance and long-term battery health can be considered together.

\subsection{Gap 5: Limited Robustness and Empirical Evaluation}

Battery models and learned policies can be affected by sensor noise,
model mismatch, operating-condition changes, and battery ageing.
Therefore, evaluation under different conditions is necessary to
determine whether a controller is robust.

\section{Research Positioning}

The identified gaps motivate the development of a Multi-Agent Battery
Management System (MA-BMS). The framework considered in this thesis
uses functional decomposition in which charging and battery-health
objectives are handled by separate learning components, while a
dedicated safety supervisor provides an additional safety layer.

The overall research direction can therefore be summarized as follows:

\begin{table}[htbp]
\centering
\caption{Research gaps and corresponding research direction.}
\label{tab:research_positioning}
\begin{tabular}{p{4.2cm}p{7.0cm}}
\toprule
\textbf{Research Gap} & \textbf{Research Direction} \\
\midrule
Ageing-unaware control &
Incorporate battery ageing information into the control process \\

Objective entanglement &
Use functional multi-agent decomposition \\

Lack of hard safety &
Use an independent safety-supervision layer \\

Separated health and charging decisions &
Integrate battery-health information into charging control \\

Limited robustness &
Evaluate the controller under different operating and uncertainty
conditions \\
\bottomrule
\end{tabular}
\end{table}

\section{Chapter Summary}

This chapter reviewed conventional battery charging, battery state
estimation, battery ageing and SoH estimation, model-based control,
reinforcement learning, multi-agent reinforcement learning, and safe
reinforcement learning.

The review shows that conventional and model-based methods provide
strong physical interpretability and constraint handling but can be
limited by fixed rules, model accuracy, and computational requirements.
Reinforcement learning provides adaptive control but commonly combines
multiple objectives through a single reward function. Multi-agent
reinforcement learning provides an opportunity to decompose these
objectives, while a separate safety mechanism can address the
limitations of purely learned control.

These observations establish the research gaps addressed by the
proposed MA-BMS framework. The following chapter presents the
methodology and system architecture used to address these gaps.